% !TEX root = ../main.tex
\lecture{1}{2026-09-14}{Day 1 Intro}

\url{https://lms.uzh.ch/auth/RepositoryEntry/17923670210/CourseNode/114458034551518/path%3D~~Aufdatierte%20Dokumente~~Vorlesungsnotizen/0}

\section{Admin stuff}
\begin{itemize}
    \item Livestream auf OLAT
    \item Podcast verfügbar
    \item erste 6 Wochen grundwissen (Mittelschule, Gymi Level)
    \item Freiwillige uebungen keine einschreibung noetig
    \item Verantwortliche fuer die Uebung werden keine Fragen zur pruefung beantworten können
    \item Fragen E-Mail an financial.accounting@business.uzh.ch 
    \item Fragen werden im Q\&A Dokument auf OLAT beantwortet.
    \item Fragen möglich bis: 9. Dezember 2026, 10:00 Uhr 
    \item \textbf{Pruefung}
    \subitem 16.12 10:00-11:00
    \subitem Messe Zueri Oerlikon
    \subitem Multiple Choice Pruefung immer nur 1 richtig
    \subitem Teil A einzelne Aussagen, B Aufgaben mit einem Fall Beispiel, C Aufgaben 2 Beispiele
\end{itemize}

\section{Grundlagen}

Um was handelt es sich im FINACC: Unternehmen

\textbf{Definition Unternehmen \& Financial Accounting}
\begin{itemize}
    \item Untersuchungsobjekt der BWL,
    \item Netzwerk von Verträgen
    \subitem Angestellte
    \subitem Kunde 
    \subitem Andere Unternehmen
    \item Rechnungswesen (Accounting)
    \subitem Abbildung der wirtschaftlichen Auswirkungen der Verträge
\end{itemize}

\subsection{Wer nutzt die Accountingdaten?}
\textbf{Entitäten \& Gründe}
\begin{itemize}
    \item \textbf{Aussen} (Abschätzung der wirtschaftlichen Lage eines Unternehmens)
        \subitem Staat
        \subitem Investoren
        \subitem Merger
    \item \textbf{Innen} (Leitung der Unternehmensstrategie und Entscheidungen)
        \subitem Aquisition
        \subitem Fabrik umzug (Kosten und Resultat)
\end{itemize}

Aus diesen zwei Sphären entstehen Rechnungswesen bereiche
\begin{itemize}
    \item \textbf{finanzielle Rechnungswesen} (AUSSEN)
        \subitem Geregelte Standards   
        \subitem OR Art. 957
    \item  \textbf{managerial Accounting} (INNEN)
        \subitem keine Standards, Intern geregelt 
\end{itemize}


\textbf{Ziele des \textit{Finanzellen Rechnungswesens}}
\begin{itemize}
    \item Infomations Asymmetrie (Innen/Aussen) abbauen durch Accdaten
    \item Guide für betriebliche Entscheidungen für INvestoren und Geschäftsführer
    \item Regelung der Accountingstandards dienen der Reduktion des Spielraums für Manipulation
    \item Transparenz und realitätstreue Berichtung
\end{itemize}

\subsection{Das Konto als Basis des Rechnungswesens}

\begin{definition}
    [Konto] Werkzeug für die systematische Abbildung der wirtschaftlichen Folgen der Unternehmensverträge
\end{definition}


\textbf{Aspekten eines standard Kontos}
\begin{itemize}
    \item Darstellungsform immer mit 2 Seiten links rechts
        \subitem 1. Soll
        \subitem 2. Haben
    \pagebreak
    \item Doppelte Buchhaltung (Tool)
    \begin{itemize}
        \item Struktur
        \begin{itemize}
            \item Kasse
            
            \item AB (Anfangsbestand) SB (Schlussbestand)
            
            \item Links Soll (Einnahmen)
            
            \item Rechts Haben (Ausgaben)
            
            \item Ein Wert pro Zeile
            
            \item Summenresultat links und rechts Spalten (inkl. AB und SB) \textbf{gleich} \footnote{Die gleichheit der Zu und Abnahmen}
            
        \end{itemize}
    \end{itemize}
\end{itemize}
            
\subsection{Buchungstatsachen}
\begin{definition}
    [Buchungstatsache] Gründe für die Veränderung der Konten. Ursprung der Veränderungen sind die Verträge
\end{definition}
\begin{itemize}
    \item Explizite Verträge
        \subitem Kaufverträge 
        \subitem Mietvertrag 
    \item Implizite Verträge
        \subitem Bewertung einer Position durch Schätzung des möglichen Werts
        \subitem Devaluation Maschinen (Konstanten Kapitals)
\end{itemize}

\textbf{Jahresberichte}
Sammlung der Buchungstatsachen zu einem Jahresbericht, damit kann die kontinuierliche finanzielle Geschichte eines Unternehmens verfolgen.

\textbf{Wie könnten die Schwankungen im Aktienkurs der ABB so start verändern?} \break
Bei der Antwort helfen die Finanzberichte
